\documentclass[oneside,openany]{dukedissertation2026}

%\documentclass[economy,twoside,bind]{dukedissertation}
% Use the second for a single-spaced copy suitable for duplex printing
% and binding.

% Other useful options (there are more options documented in Chapter 2):
%  * draft -- don't actually include images, print a black bar on overful
%             hboxes
%  * MS    -- Format for a Master's Thesis.  No UMI abstract page, some 
%             textual changes to title page.  


% Useful packages for dissertation writing:
% fontenc is not needed because the class uses fontspec/unicode-math (compile with LuaLaTeX or XeLaTeX)
\usepackage{titlesec}
\usepackage{bm}
\usepackage{graphicx}

\usepackage[numbers,square,sort&compress]{natbib}


\usepackage[title,titletoc]{appendix}
\usepackage{apptools}
%\usepackage{mathptmx} % Uncomment this to use Times instead of Palatino
\usepackage{color}
\usepackage{bm}

\usepackage{multirow}
\usepackage{setspace}


\usepackage{indentfirst}
\usepackage{enumitem}
\usepackage{xfrac}
\usepackage{mathtools}
\usepackage{caption}
\usepackage{subcaption}
\usepackage{tabularx}
\usepackage{booktabs}
\usepackage{siunitx}
%\usepackage{lscape}
% Use the pdflscape package if you want the page to rotate
\usepackage{pdflscape}
%\usepackage{longtable}
\usepackage{listings}
\usepackage{xcolor}
\usepackage{pgfplots}
%\usepackage{acronym}
%\usepackage{showframe}
\usepackage{makecell}
\usepackage[super]{nth}
\usepackage{xr}

\setlist[itemize]{noitemsep,nolistsep}
\setlist[enumerate]{noitemsep,nolistsep}
\sisetup{per-mode = symbol}
\sisetup{uncertainty-mode = separate}
\sisetup{
	output-open-uncertainty = [,
	output-close-uncertainty = ],
	uncertainty-separator = \,
}

%\usepackage{cite}  % If you include this, hyperlink cites will
                     % break.  It's nice to use this package if your bibstyle
							% sorts entries by order-of-use, rather than
							% alphabetically (as plain does).
							
%Theorem, Lemma, etc. environments
\newtheorem{theorem}{Theorem}%[section]
\newtheorem{lemma}[theorem]{Lemma}
\newtheorem{proposition}[theorem]{Proposition}
\newtheorem{corollary}[theorem]{Corollary}
\newtheorem{result}[theorem]{Result}

% Personal commands and abbreviations.
%Define and personal commands here

%Graphics Path to find your pictures
\graphicspath{{./Pictures/}{./Chapter1}} %Add multiple paths if there are images in multiple spots
% Change CaptionSetup justification to one of: raggedright, centering, or justified.
\captionsetup{justification=centering, singlelinecheck=false} 


%-----------------------------------------------------------------------------%
% PREAMBLE 
%-----------------------------------------------------------------------------%
\author{Huarui Zhou}
\title{[Your Title Here Even If It Is Very Very Very Very Very Very Long. Truly It Is A Miraculously Long Title For A Dissertation]}
\supervisor{Steven B. Haase}
\department{Biology} % Appears as Department of \department
\program{Cell and Molecular Biology}
% Declare dissertation subject used on UMI abstract page.  List of
% categories: http://dissertations.umi.com/duke/subject_categories.html
%\subject{[Your Subject Here]}

\date{2027} % Anything but the year is ignored.

% Copyright text.  If undefined, default is 'All rights reserved'
% (Example sets the text to a hyperlinked Creative Commons Licence)
\copyrighttext{All rights reserved except the rights granted by the \href{http://creativecommons.org/licenses/by-nc/3.0/us/}{Creative Commons Attribution-Noncommercial Licence}}

% Committee Members other than supervisor.  No more than five beyond the
% supervisor allowed.Do not include titles (e.g. `Dr.')
\defensedate{May 20, 2027}
\member{[Committee Member 1]}
\member{[Committee Member 2]}
\member{[Committee Member 3]}
\member{[Committee Member 4]}
%-----------------------------------------------------------------------------%

%-----------------------------------------------------------------------------%
% HYPERREF: plain black hypertext references for ref's and cite's.
%-----------------------------------------------------------------------------%
\usepackage[pdfusetitle, plainpages=false, letterpaper, bookmarks, bookmarksnumbered, colorlinks, linkcolor=black, citecolor=black, filecolor=black, urlcolor=black,linktoc=none]{hyperref}
\usepackage[capitalise]{cleveref}
\creflabelformat{equation}{#2#1#3}
\makeatletter

\newcommand\getfontsizeofheading{The current font size is: \fontname\font \f@size pt}
% print given font size in point and in its "font size" (i.e. larger or smaller)
\newcommand\getfontsizeforfontsizeclass[1]{{\string #1 is printed as #1  \fontname\font \f@size pt\par}}
% print the entire string in the size of the fontsize (i.e. larger or smaller)
\newcommand\printfontsizeforfontsizeclass[1]{{#1 \string #1 is printed in  \fontname\font \f@size pt\par}}
\makeatother

%\AtBeginBibliography{\vspace*{\baselineskip}}
%\setlength{\bibitemsep}{\baselineskip}

%\titleformat{\chapter}
%	{\Large\bfseries\usefont{T1}{phv}{b}{n}\fontsize{16pt}{\baselineskip}} %Format
%	{\thechapter.} %Label
%	{\fontdimen2\font\space} %Sep
%	{\normalbaselines} % Before 

%\renewcommand{\thesection}{\arabic{chapter}.\arabic{section}}
%\renewcommand{\thesubsection}{\arabic{chapter}.\arabic{section}.\arabic{subsection}}
%\renewcommand{\thesubsubsection}{\arabic{chapter}.\arabic{section}.\arabic{subsection}.\arabic{subsubsection}}

\titleformat{\section}[block]{\normalfont\bfseries\fontsize{14pt}{16pt}\selectfont}{\thesection}{.2em}{\normalbaselines}
\titleformat{\subsection}[block]{\normalfont\bfseries\fontsize{13pt}{15pt}\selectfont}{\thesubsection}{.2em}{\normalbaselines}
\titleformat{\subsubsection}[block]{\normalfont\bfseries\fontsize{11pt}{13pt}\selectfont}{\thesubsubsection}{.2em}{\normalbaselines}

\titleformat{\paragraph}[runin]{\normalfont\bfseries\fontsize{11pt}{13pt}\selectfont}{\theparagraph}{.2em}{}
\titleformat{\subparagraph}[runin]{\normalfont\bfseries\fontsize{11pt}{13pt}\selectfont}{\thesubparagraph}{.2em}{}
  
  
\titlespacing\section{0pt}{14pt plus 0pt minus 0pt}{6pt plus 0pt minus 0pt}
\titlespacing\subsection{0pt}{14pt plus 0pt minus 0pt}{6pt plus 0pt minus 0pt}
\titlespacing\subsubsection{0pt}{14pt plus 0pt minus 0pt}{6pt plus 0pt minus 0pt}
\titlespacing\paragraph{0pt}{6pt plus 0pt minus 0pt}{6pt plus 0pt minus 0pt}
\titlespacing\subparagraph{0pt}{6pt plus 0pt minus 0pt}{6pt plus 0pt minus 0pt}

\setcounter{secnumdepth}{4}

\newcolumntype{L}[1]{>{\raggedright\let\newline\\\arraybackslash\hspace{0pt}}b{#1}}
\newcolumntype{C}[1]{>{\centering\let\newline\\\arraybackslash\hspace{0pt}}b{#1}}
\newcolumntype{R}[1]{>{\raggedleft\let\newline\\\arraybackslash\hspace{0pt}}b{#1}}

\usepackage{subfiles}

\begin{document}

%-----------------------------------------------------------------------------%
% TITLE PAGE -- provides UMI abstract title page & copyright if appropriate
%-----------------------------------------------------------------------------%
\maketitle

%-----------------------------------------------------------------------------%
% ABSTRACT -- included file should start with '\abstract'.
%-----------------------------------------------------------------------------%
\include{./Abstract/abstract}
%-----------------------------------------------------------------------------%
% DEDICATION -- OPTIONAL.  Put the text inside the braces.
%               (Long 'dedications' probably belong in the acknowledgements)
%-----------------------------------------------------------------------------%
\dedication{If you want to dedicate your thesis to anyone do so here.}

%-----------------------------------------------------------------------------%
% FRONTMATTER -- ToC is required, LoT and LoF are required if you have any
% tables or figures, respectively. List of Abbreviations and Symbols is 
% optional.
%-----------------------------------------------------------------------------%
\tableofcontents % Automatically generated

\listoftables	% If you have any tables, automatically generated

\listoffigures	% If you have any figures, automatically generated

%\include{./Abbreviations/listofabbr} % List of Abbreviations. Start file with '\abbreviations'

%-----------------------------------------------------------------------------%
% ACKNOWLEDGEMENTS -- included file should start with '\acknowledgements'
%-----------------------------------------------------------------------------%
\include{./Acknowledgements/acknowledgements}

%==============================================================================
%-----------------------------------------------------------------------------%
%
% MAIN BODY OF PAPER
%
%
%-----------------------------------------------------------------------------%
\include{Chapter1/introduction}
\chapter{Your title of Chapter 2}
\label{Chapter2} % A regular chapter, starts with '\chapter{Title}'
\input{Chapter3/chRegularLook}


%==============================================================================

%-----------------------------------------------------------------------------%
% APPENDICES -- OPTIONAL. These are just chapters enumerated by Appendix A,
%                Appendix B, Appendix C...
%-----------------------------------------------------------------------------%

\begin{appendices}
	\titleformat{\chapter}[block]
	{\normalfont\bfseries\fontsize{16pt}{20pt}\selectfont}
	{Appendix \thechapter.}{.2em}{\normalbaselines}
	
	\include{./Appendix1/chLorem3} % Start with '\chapter{Title}'
\end{appendices}


%You can always add more appendices here if you want

%-----------------------------------------------------------------------------%
% BIBLIOGRAPHY -- uncomment \nocite{*} to include items in 'mybib.bib' file
% that aren't cited in the text.  Change the style to match your
% discipline's standards.  Of course, if your bibliography file isn't called
% 'mybib.bib' you might want to change that here too :)
%-----------------------------------------------------------------------------%
%\nocite{*} %- if you use this it will put EVERYTHING in your .bib file into the references even if you don't cite it in the text
\clearpage
\normalbaselines

\renewcommand{\bibname}{References}
\addcontentsline{toc}{chapter}{References}

%\nocite{*}
\bibliographystyle{naturemag}
\bibliography{Bibliography/mybib}
%-----------------------------------------------------------------------------%

%-----------------------------------------------------------------------------%
% BIOGRAPHY -- Start file with '\biography'.  Mandatory for Ph.D.
%-----------------------------------------------------------------------------%
\include{Biography/biography}

%-----------------------------------------------------------------------------
% You're done :)
\end{document}